Šīs nodaļas eksperimentālais novērtējums pārbauda promocijas darba realizācijas, izmantojot \ref{sec:impl}. nodaļā aplūkoto
realizāciju rezultātus. Tā kā piedāvātās realizācijas pievēršas dažādiem neiromorfās skaitļošanas aspektiem, pārbaude
sadalīta trijās saistītās daļās.

Pirmajā daļā pārbaudīts darbā \cite{lielamurs2026} piedāvātais \Gls{fpga} paātrinātājs. Tā realizācijas veiktspēja mērīta ar
attēlu klasifikācijas veiktspējas testiem, analizējot inferences latentumu, caurlaidību, resursu izmantojumu, kvantēšanas ietekmi
un novērtēto enerģijas patēriņu.

Otrajā daļā pārbaudīta piedāvātā telpiski temporālā objektu detektēšanas metode LiDAR punktu mākoņiem \cite{lielamurs2024}. Tajā uzmanība pievērsta detektēšanas precizitātei un salīdzināta dziļo konvolucionālo impulsu neironu tīklu un tradicionālo mākslīgo neironu tīklu veiktspēja \textit{KITTI} datu kopā.

Trešajā daļā pārbaudīta piedāvātā detektēšanas ietvara piemērotība lidojuma laika sensoru datiem \cite{lielamurs2025}. Eksperimentālie rezultāti ar automobiļu \Gls{tof} kameru parāda laikā kodētu \Gls{snn} piemērotību praktiskiem 3D uztveres uzdevumiem.
